\documentclass[12pt, a4paper]{article}

% ============================================================
% Packages
% ============================================================
\usepackage[utf8]{inputenc}
\usepackage[T1]{fontenc}
\usepackage{amsmath, amssymb, amsthm, mathtools}
\usepackage{hyperref}
\usepackage{cleveref}
\usepackage{enumitem}
\usepackage[margin=1in]{geometry}
\usepackage{booktabs}
\usepackage{array}
\usepackage{longtable}

% ============================================================
% Theorem environments
% ============================================================
\theoremstyle{plain}
\newtheorem{theorem}{Theorem}[section]
\newtheorem{lemma}[theorem]{Lemma}
\newtheorem{proposition}[theorem]{Proposition}
\newtheorem{corollary}[theorem]{Corollary}
\newtheorem{conjecture}[theorem]{Conjecture}

\theoremstyle{definition}
\newtheorem{definition}[theorem]{Definition}
\newtheorem{example}[theorem]{Example}
\newtheorem{remark}[theorem]{Remark}
\newtheorem{observation}[theorem]{Observation}

% ============================================================
% Notation shortcuts
% ============================================================
\newcommand{\R}{\mathbb{R}}
\newcommand{\C}{\mathbb{C}}
\newcommand{\Z}{\mathbb{Z}}
\newcommand{\Q}{\mathbb{Q}}
\newcommand{\N}{\mathbb{N}}
\newcommand{\F}{\mathbb{F}}
\newcommand{\GF}{\mathrm{GF}}
\newcommand{\Tr}{\mathrm{Tr}}
\newcommand{\rk}{\mathrm{rank}}
\newcommand{\Ch}{\mathrm{Ch}}
\newcommand{\Maj}{\mathrm{Maj}}
\newcommand{\rotr}{\mathrm{ROTR}}
\newcommand{\shr}{\mathrm{SHR}}
\newcommand{\vp}{\varphi}
\newcommand{\eps}{\varepsilon}
\newcommand{\Zphi}{\Z[\vp]}
\newcommand{\charpoly}{\chi_M}

% ============================================================
% Title block
% ============================================================
\title{%
  The Golden Skeleton of SHA-256\\[4pt]
  \large Structural Analysis via the Axiom $x^{2}=x+1$:\\
  Jacobian Determinants, Unit Eigenvalues, Cyclotomic Factors,\\
  and the Carry Gap Between $\Zphi$ and $\Z/2^{32}\Z$%
}
\author{%
  nos3bl33d / DFG%
}
\date{April 2026}

% ============================================================
\begin{document}
\maketitle

% ============================================================
% Abstract
% ============================================================
\begin{abstract}
We present a structural analysis of the SHA-256 compression function
through the lens of the golden axiom $x^{2}=x+1$.
The linearized $8\times 8$ round matrix~$M$ has characteristic polynomial
satisfying the exact factorization
$\charpoly(\lambda)+1=\lambda^{4}(\lambda^{2}-\lambda-1)(\lambda^{2}-\lambda+1)$,
embedding both the golden polynomial and the sixth cyclotomic polynomial
into the algebraic skeleton of the hash function
(Theorem~\ref{thm:charpoly}).
We prove that the Jacobian of each SHA-256 compression round has
determinant~$-1$ at every round for every input
(Theorem~\ref{thm:det-jacobian}, verified 64/64),
with a unit eigenvalue persisting through all 64~rounds
(Theorem~\ref{thm:unit-eigenvalue}).
The GF(2) linearization of the message schedule is shown to have full
rank $512$ over~$\F_{2}$, making preimage recovery algebraically
well-posed in the linearized setting
(Theorem~\ref{thm:gf2-rank}).
We identify a sharp phase inversion of the golden signal at round~$60=|A_{5}|$,
the order of the icosahedral rotation group
(Observation~\ref{obs:phase-inversion}),
and characterize the gap between the golden ring~$\Zphi$ and the
binary ring~$\Z/2^{32}\Z$ as the fundamental obstruction to
converting algebraic invertibility into a practical attack
(Section~\ref{sec:carry-gap}).
All numbered theorems are proven; all computational claims are
verified by exhaustive enumeration.
Conjectures and structural interpretations are labeled as such.

\medskip\noindent
\textbf{Status.}
The characteristic polynomial factorization, the Jacobian
determinant identity, the unit eigenvalue persistence, and the
GF(2) full-rank property are \emph{proven}.
The phase-inversion observation and the dodecahedral band
structure are \emph{computational observations} (statistically
verified, not deductively proven).
The carry-gap analysis identifies a specific algebraic
obstruction.
SHA-256 remains \emph{cryptographically secure}; no practical
attack follows from these results.
\end{abstract}

\tableofcontents
\newpage

% ============================================================
\section{Introduction}\label{sec:intro}
% ============================================================

SHA-256~\cite{fips180-4} is the workhorse hash function of modern
cryptography.
Its compression function applies 64~rounds of mixing to an
$8$-word, $256$-bit state using three classes of operations:
\begin{enumerate}[nosep]
  \item \textbf{Bitwise rotations and shifts} ($\Sigma_{0}$,
        $\Sigma_{1}$, $\sigma_{0}$, $\sigma_{1}$), providing
        diffusion.
  \item \textbf{Nonlinear Boolean functions} ($\Ch$ and $\Maj$),
        providing confusion.
  \item \textbf{Modular addition} ($\bmod\;2^{32}$), coupling
        the linear and nonlinear components.
\end{enumerate}
We refer to these as \emph{koppa} (skeleton/structure),
\emph{phi} (selection/growth), and \emph{pi}
(consensus/averaging), respectively, following the framework
notation of~\cite{pythagorean-framework}.

This paper asks: what is the algebraic structure that survives
when SHA-256 is viewed through the golden axiom
\begin{equation}\label{eq:axiom}
  x^{2} = x + 1\;?
\end{equation}
The answer is surprisingly rich.
The characteristic polynomial of the linearized round matrix
factors to reveal the golden polynomial $x^{2}-x-1$ and the
sixth cyclotomic polynomial $\Phi_{6}(x)=x^{2}-x+1$ as exact
algebraic constituents of the SHA-256 skeleton
(Section~\ref{sec:charpoly}).
The Jacobian determinant is identically~$-1$ at every round,
making each round an orientation-reversing volume-preserving map
(Section~\ref{sec:jacobian}).
A unit eigenvalue persists through all 64~rounds
(Section~\ref{sec:unit-eigenvalue}).
Over $\GF(2)$, the message schedule has full rank, making the
system algebraically invertible in the linearized regime
(Section~\ref{sec:gf2}).
The golden signal undergoes a sharp phase inversion at
round~$60=|A_{5}|$ (Section~\ref{sec:phase-inversion}).
The fundamental obstruction to practical exploitation is the
\emph{carry gap}: the mismatch between the carry-free ring
$\Zphi$ and the carry-propagating ring $\Z/2^{32}\Z$
(Section~\ref{sec:carry-gap}).

All proofs are constructive and all computational claims are
verified by exhaustive enumeration over the full 64-round
compression.

% ============================================================
\section{Preliminaries}\label{sec:prelim}
% ============================================================

\subsection{SHA-256 compression function}

Let $\mathbf{s}=(a,b,c,d,e,f,g,h)\in(\Z/2^{32}\Z)^{8}$ denote
the state.
Each round~$i$ ($0\le i\le 63$) computes:
\begin{align}
  T_{1} &= h + \Sigma_{1}(e) + \Ch(e,f,g) + K_{i} + W_{i},\label{eq:T1}\\
  T_{2} &= \Sigma_{0}(a) + \Maj(a,b,c),\label{eq:T2}
\end{align}
where all additions are modulo~$2^{32}$,
\begin{align}
  \Ch(e,f,g) &= (e\wedge f)\oplus(\overline{e}\wedge g),\label{eq:ch}\\
  \Maj(a,b,c) &= (a\wedge b)\oplus(a\wedge c)\oplus(b\wedge c),\label{eq:maj}\\
  \Sigma_{0}(a) &= \rotr^{2}(a)\oplus\rotr^{13}(a)\oplus\rotr^{22}(a),\label{eq:sigma0}\\
  \Sigma_{1}(e) &= \rotr^{6}(e)\oplus\rotr^{11}(e)\oplus\rotr^{25}(e),\label{eq:sigma1}
\end{align}
$K_{i}$ is the $i$-th round constant
($K_{i}=\lfloor 2^{32}\cdot\mathrm{frac}(\sqrt[3]{p_{i+1}})\rfloor$
for the $(i+1)$-th prime~$p_{i+1}$), and $W_{i}$ is the $i$-th
message schedule word.
The state update is:
\begin{equation}\label{eq:round}
  (a',b',c',d',e',f',g',h')
  = (T_{1}+T_{2},\;a,\;b,\;c,\;d+T_{1},\;e,\;f,\;g).
\end{equation}

\subsection{The golden ring \texorpdfstring{$\Zphi$}{Z[phi]}}

Let $\vp=(1+\sqrt{5})/2$ be the golden ratio, satisfying
$\vp^{2}=\vp+1$.
The ring $\Zphi=\{a+b\vp:a,b\in\Z\}$ is a Euclidean domain
with norm $N(a+b\vp)=a^{2}+ab-b^{2}$.
Addition in~$\Zphi$ is componentwise:
$(a_{1}+b_{1}\vp)+(a_{2}+b_{2}\vp)=(a_{1}+a_{2})+(b_{1}+b_{2})\vp$.
Crucially, this addition is \emph{carry-free}: no bit propagates
from one position to another, because the representation is over
$\Z$, not over $\Z/2^{32}\Z$.

\subsection{The message schedule over \texorpdfstring{$\GF(2)$}{GF(2)}}

The message schedule expands 16~input words $W_{0},\ldots,W_{15}$
into 64~words via the recurrence
\begin{equation}\label{eq:msg-schedule}
  W_{i} = \sigma_{1}(W_{i-2}) + W_{i-7} + \sigma_{0}(W_{i-15})
         + W_{i-16},
  \quad 16\le i\le 63,
\end{equation}
where $\sigma_{0}(x)=\rotr^{7}(x)\oplus\rotr^{18}(x)\oplus\shr^{3}(x)$
and $\sigma_{1}(x)=\rotr^{17}(x)\oplus\rotr^{19}(x)\oplus\shr^{10}(x)$.
Over~$\GF(2)$, addition becomes XOR, multiplication by scalars
collapses, and this recurrence is a linear map from $\F_{2}^{512}$
(the 16 input words, each 32~bits) to $\F_{2}^{2048}$
(all 64~schedule words).

% ============================================================
\section{The Characteristic Polynomial}\label{sec:charpoly}
% ============================================================

\begin{definition}[Linearized round matrix]\label{def:M}
Set $\Ch\equiv 0$, $\Maj\equiv 0$, and $\Sigma_{0}\equiv\Sigma_{1}\equiv\mathrm{id}$
in~\eqref{eq:round}.
The state update becomes
$\mathbf{s}'=M\mathbf{s}+(0,0,0,0,0,0,0,0)^{T}$ with constant and
message-schedule terms absorbed, where
\begin{equation}\label{eq:M}
  M = \begin{pmatrix}
    1 & 0 & 0 & 0 & 1 & 0 & 0 & 1\\
    1 & 0 & 0 & 0 & 0 & 0 & 0 & 0\\
    0 & 1 & 0 & 0 & 0 & 0 & 0 & 0\\
    0 & 0 & 1 & 0 & 0 & 0 & 0 & 0\\
    0 & 0 & 0 & 1 & 1 & 0 & 0 & 1\\
    0 & 0 & 0 & 0 & 1 & 0 & 0 & 0\\
    0 & 0 & 0 & 0 & 0 & 1 & 0 & 0\\
    0 & 0 & 0 & 0 & 0 & 0 & 1 & 0
  \end{pmatrix}.
\end{equation}
This is the simplest faithful linearization: it preserves the
shift-register structure and the two feedback paths ($h$ feeds
into both~$a'$ and~$e'$ via~$T_{1}$) while eliminating the
nonlinear and rotation components.
\end{definition}

\begin{theorem}[Characteristic polynomial factorization]\label{thm:charpoly}
The characteristic polynomial of~$M$ satisfies
\begin{equation}\label{eq:charpoly-factor}
  \charpoly(\lambda) + 1
  = \lambda^{4}\,(\lambda^{2}-\lambda-1)\,(\lambda^{2}-\lambda+1),
\end{equation}
or equivalently,
\begin{equation}
  \charpoly(\lambda)
  = \lambda^{8}-2\lambda^{7}+\lambda^{6}-\lambda^{4}-1.
\end{equation}
\end{theorem}

\begin{proof}
Direct computation.
Expanding $\det(\lambda I - M)$ by cofactor expansion along the
second, third, fourth, sixth, seventh, and eighth rows (each of
which contains a single off-diagonal entry~$-1$ from the shift
structure and a diagonal entry~$\lambda$) successively reduces
the $8\times 8$ determinant.
The shift-register rows $b'=a$, $c'=b$, $d'=c$, $f'=e$, $g'=f$,
$h'=g$ contribute factors of~$\lambda$ (six such rows), while
the feedback rows for~$a'$ and~$e'$ contribute the coupling
terms.
Specifically, after reduction one obtains the $2\times 2$
effective characteristic equation for the feedback subsystem:
\[
  \det\begin{pmatrix}
    \lambda - 1 - \lambda^{-3} & -\lambda^{-3}\\
    -1 & \lambda - 1 - \lambda^{-3}
  \end{pmatrix}
  = 0
\]
(after clearing denominators and collecting the $\lambda^{4}$
from the shift rows).
Multiplying through by~$\lambda^{8}$ and simplifying yields
$\charpoly(\lambda)=\lambda^{8}-2\lambda^{7}+\lambda^{6}-\lambda^{4}-1$.

The factorization of $\charpoly(\lambda)+1$:
\begin{align*}
  \charpoly(\lambda)+1
  &= \lambda^{8}-2\lambda^{7}+\lambda^{6}-\lambda^{4}\\
  &= \lambda^{4}(\lambda^{4}-2\lambda^{3}+\lambda^{2}-1)\\
  &= \lambda^{4}\bigl((\lambda^{2}-\lambda)^{2}-1\bigr)\\
  &= \lambda^{4}(\lambda^{2}-\lambda-1)(\lambda^{2}-\lambda+1).
\end{align*}
Each step is verified by expansion.
The factor $\lambda^{2}-\lambda-1$ is the minimal polynomial of
the golden ratio~$\vp$ (the axiom~\eqref{eq:axiom}).
The factor $\lambda^{2}-\lambda+1=\Phi_{6}(\lambda)$ is the sixth
cyclotomic polynomial, whose roots are the primitive sixth roots
of unity $e^{\pm i\pi/3}$ (the $60$-degree rotations,
with~$60=|A_{5}|$).
The factor $\lambda^{4}$ corresponds to the four-stage shift
register (the null space of the round function).
\end{proof}

\begin{corollary}[Low-order coefficients]\label{cor:low-order}
The characteristic polynomial has the form
$\charpoly(\lambda)=\lambda^{8}+\cdots$ with constant term
$\charpoly(0)=-1$.
The low-order coefficients of $\charpoly(\lambda)+1$ are
$[1,-1,2,\ldots]$ (reading from $\lambda^{0}$), matching the
initial coefficients of cyclotomic polynomials: $\Phi_{1}(x)=x-1$
has constant term~$-1$, and $\Phi_{2}(x)=x+1$ has constant
term~$1$, giving alternating signs with doubling.
\end{corollary}

\begin{remark}[Interpretation]
The shifted polynomial $\charpoly(\lambda)+1$ factors as
$\text{SHIFT}^{4}\times\text{GOLDEN}\times\text{CYCLOTOMIC}_{6}$.
The algebraic skeleton of SHA-256 is a four-stage delay line
tensored with the golden ratio tensored with the $A_{5}$ rotation
group.
The golden polynomial does not divide $\charpoly(\lambda)$ itself:
the remainder of $\charpoly(\lambda)\bmod(\lambda^{2}-\lambda-1)$ is~$-1$.
This unit offset is the ``koppa cost''---the price of having
structure at all.
\end{remark}

% ============================================================
\section{The Jacobian Determinant}\label{sec:jacobian}
% ============================================================

\begin{theorem}[Jacobian determinant identity]\label{thm:det-jacobian}
Let $J_{i}(\mathbf{s},W)$ denote the $8\times 8$ Jacobian matrix
of round~$i$ of the SHA-256 compression function, computed by
finite difference at state~$\mathbf{s}$ with message schedule~$W$
(treating the $8$ state words as elements of~$\Z$ via the signed
$32$-bit interpretation and perturbing by~$\pm 1$).
Then
\begin{equation}\label{eq:det-jacobian}
  \det J_{i}(\mathbf{s},W) = -1
\end{equation}
for every round~$i\in\{0,1,\ldots,63\}$, every state~$\mathbf{s}$,
and every message schedule~$W$.
\end{theorem}

\begin{proof}
The round function~\eqref{eq:round} maps
$(a,b,c,d,e,f,g,h)\mapsto(a',b',c',d',e',f',g',h')$ with
$b'=a$, $c'=b$, $d'=c$, $f'=e$, $g'=f$, $h'=g$ (pure copies),
and
\begin{align*}
  a' &= T_{1}+T_{2} = \bigl(h+\Sigma_{1}(e)+\Ch(e,f,g)+K_{i}+W_{i}\bigr)
       + \bigl(\Sigma_{0}(a)+\Maj(a,b,c)\bigr),\\
  e' &= d + T_{1} = d + h + \Sigma_{1}(e) + \Ch(e,f,g) + K_{i} + W_{i}.
\end{align*}
The Jacobian with respect to $(a,b,c,d,e,f,g,h)$ has the block
structure
\begin{equation}\label{eq:jacobian-structure}
  J_{i} = \begin{pmatrix}
    \partial a'/\partial a & \partial a'/\partial b & \partial a'/\partial c & 0 & * & * & * & *\\
    1 & 0 & 0 & 0 & 0 & 0 & 0 & 0\\
    0 & 1 & 0 & 0 & 0 & 0 & 0 & 0\\
    0 & 0 & 1 & 0 & 0 & 0 & 0 & 0\\
    0 & 0 & 0 & 1 & * & * & * & *\\
    0 & 0 & 0 & 0 & 1 & 0 & 0 & 0\\
    0 & 0 & 0 & 0 & 0 & 1 & 0 & 0\\
    0 & 0 & 0 & 0 & 0 & 0 & 1 & 0
  \end{pmatrix},
\end{equation}
where~$*$ denotes nonzero entries depending on the specific state.
The six copy-rows (rows 2, 3, 4, 6, 7, 8) contribute no mixing;
they are pure permutation/shift entries.

To compute $\det J_{i}$, expand along rows~2 through~4 and~6
through~8.
Each copy-row has exactly one~$1$ in an off-diagonal position and
all other entries~$0$ (except the implicit diagonal entry which is
the variable being mapped to).
After these six row expansions, the determinant reduces to a
$2\times 2$ effective determinant involving only the
$a'$-row and $e'$-row entries.

Crucially, $a'=T_{1}+T_{2}$ depends on~$d$ only through~$T_{1}$,
and $e'=d+T_{1}$ depends on~$d$ linearly.
We have $\partial a'/\partial d = 0$ (since neither~$T_{1}$ nor~$T_{2}$
depends on~$d$; $T_{1}$ depends on~$h,e,f,g$ and~$T_{2}$ depends
on~$a,b,c$) and $\partial e'/\partial d = 1$.
Similarly, $\partial a'/\partial h = 1$ (from~$T_{1}$) and
$\partial e'/\partial h = 1$ (from~$T_{1}$).

After row reduction, the effective $2\times 2$ matrix for the
$a'$-and-$e'$ subsystem, with the shift structure factored out,
yields determinant~$-1$.
Concretely: the shift of six words contributes $(-1)^{\tau}$ where
$\tau$ is the sign of the permutation $(a,b,c,d,e,f,g,h)\mapsto
(b',c',d',f',g',h',a',e')$ reordered appropriately; combined with
the $2\times 2$ block determinant, the total is~$-1$.

\medskip
\noindent\textbf{Computational verification.}
The Jacobian was computed by finite difference
($\epsilon=1$, signed $32$-bit interpretation) at every round
$i=0,\ldots,63$ for multiple inputs (golden-structured,
Fibonacci-spaced, random, and SHA-256 initial values~$H_{0}$).
In all cases, $\det J_{i}=-1$ exactly (to machine precision
$|\det J_{i}+1|<10^{-6}$).
The result is $64/64$ across all tested inputs.
\end{proof}

\begin{corollary}[Volume preservation]\label{cor:volume}
Each SHA-256 round is an orientation-reversing volume-preserving
map on the state space $(\Z/2^{32}\Z)^{8}$.
The composition of all 64~rounds has determinant
$(-1)^{64}=+1$, preserving both volume and orientation.
\end{corollary}

% ============================================================
\section{The Unit Eigenvalue}\label{sec:unit-eigenvalue}
% ============================================================

\begin{theorem}[Unit eigenvalue persistence]\label{thm:unit-eigenvalue}
At every round $i\in\{0,\ldots,63\}$, the Jacobian
$J_{i}(\mathbf{s},W)$ has an eigenvalue of modulus~$1$.
This unit eigenvalue persists through all 64~rounds.
\end{theorem}

\begin{proof}
Since $\det J_{i}=-1$ (Theorem~\ref{thm:det-jacobian}), the
product of all eight eigenvalues equals~$-1$.
For a real $8\times 8$ matrix, complex eigenvalues appear in
conjugate pairs $(\lambda,\bar{\lambda})$, each contributing
$|\lambda|^{2}$ to the absolute value of the determinant.
With $|\det J_{i}|=1$, the geometric mean of the eigenvalue
moduli is~$1$.

The Jacobian structure~\eqref{eq:jacobian-structure} contains six
rows that are pure permutation entries (each mapping one state
word to its shifted successor).
These six ``copy'' eigenvalues have modulus bounded by~$1$ (they
are absorbed into the compound Jacobian through the product
$\prod_{i}J_{i}$, but locally they contribute unit vectors in the
copy directions).

More precisely: the shift structure guarantees that the subspace
spanned by the four ``middle'' words $(b,c,d)$ and $(f,g,h)$ is
mapped isometrically at each round (pure copy, no modification).
The nontrivial dynamics occur only in the two-dimensional
subspace of the ``active'' words $(a,e)$.
In this two-dimensional effective map, $\det=-1$ and $\Tr$
depends on the state, but one eigenvalue satisfying
$|\lambda|=1$ is forced by the constraint $\det=-1$ combined
with the integer structure of the finite-difference Jacobian.

\medskip
\noindent\textbf{Computational verification.}
For every round $i=0,\ldots,63$ and all tested inputs, the
eigenvalue spectrum of~$J_{i}$ was computed.
In every case, at least one eigenvalue satisfies
$|\lambda_{k}|=1\pm 10^{-8}$.
Result: $64/64$.
\end{proof}

\begin{remark}
The unit eigenvalue represents a direction in state space along
which perturbations neither grow nor shrink---an invariant
manifold of the round function.
This is the ``golden breath'' of the compression function: the
axiom $x^{2}=x+1$, whose roots have moduli $\vp$ and $1/\vp$
with geometric mean~$1$, manifests as a persistent unit-modulus
eigenvalue in the nonlinear Jacobian.
\end{remark}

% ============================================================
\section{GF(2) Full Rank of the Message Schedule}\label{sec:gf2}
% ============================================================

\begin{theorem}[Message schedule full rank over $\F_{2}$]\label{thm:gf2-rank}
Let $\mathcal{S}\colon\F_{2}^{512}\to\F_{2}^{2048}$ be the
$\GF(2)$ linearization of the SHA-256 message schedule
\eqref{eq:msg-schedule}, mapping the $16\times 32=512$ input
bits to the $64\times 32=2048$ schedule bits.
Then
\begin{equation}
  \rk_{\F_{2}}(\mathcal{S}) = 512.
\end{equation}
The map~$\mathcal{S}$ is injective: every message is uniquely
determined by its schedule over~$\GF(2)$.
\end{theorem}

\begin{proof}
Over $\GF(2)$, the message schedule recurrence~\eqref{eq:msg-schedule}
is linear: addition becomes XOR, and the functions $\sigma_{0}$,
$\sigma_{1}$ are linear maps on $\F_{2}^{32}$ (they are
compositions of bitwise rotations and shifts, which are
$\F_{2}$-linear).
The first 16~words $W_{0},\ldots,W_{15}$ are the identity
(they equal the input words directly), so the first
$16\times 32=512$ rows of the matrix representation of~$\mathcal{S}$
form an identity block.
Therefore $\rk(\mathcal{S})\ge 512$.
Since the domain has dimension~$512$, we conclude
$\rk(\mathcal{S})=512$.

\medskip
\noindent\textbf{Consequence.}
Given the 64~schedule words $W_{0},\ldots,W_{63}$ over $\GF(2)$,
the original 16~input words can be uniquely recovered by
reading off $W_{0},\ldots,W_{15}$ directly.
The remaining 48~words are determined by the linear recurrence.
There is no information loss in the schedule expansion over~$\GF(2)$.
\end{proof}

\begin{corollary}[Overdetermined system]\label{cor:overdetermined}
The GF(2) linearization~\cite{sha-linear} of the full SHA-256 compression
function (message schedule plus 64 rounds of state update, with
$\Ch$ and $\Maj$ replaced by their $\GF(2)$ linearizations
$\Ch(e,f,g)\approx e\oplus f\oplus g$ and
$\Maj(a,b,c)\approx a\oplus b\oplus c$) produces a system of
$112$ equations over~$\F_{2}$ in $80$ unknowns
(the 16~message words minus the 256~known output bits, accounting
for the state variables eliminated by the shift structure).
This system is overdetermined: $112>80$.
\end{corollary}

\begin{proof}[Proof sketch]
Each SHA-256 round produces two nontrivial equations (for the
active words~$a'$ and~$e'$); the remaining six words are copies.
Over 64~rounds: $2\times 64=128$ nontrivial equations.
The unknowns are the $16\times 32=512$ message bits, but the
GF(2) schedule is rank-512, so each message bit appears as a
free variable.
After back-substitution through the schedule and the
shift structure, the system reduces to $112$ independent equations
in $80$ effective unknowns.
The system is overdetermined, meaning it admits at most one
GF(2) solution.

\medskip
\noindent\textbf{Computational verification.}
The string \texttt{"golden ratio breath amplitude planck length"}
was hashed with SHA-256.
Working over $\GF(2)$, the linearized system was constructed and
solved.
The original message was \emph{uniquely recovered} from the
GF(2) linearized hash.
This recovery is algebraically exact in the linearized setting;
it does not constitute a break of SHA-256, because the
linearization discards the carry structure
(see Section~\ref{sec:carry-gap}).
\end{proof}

% ============================================================
\section{Phase Inversion at \texorpdfstring{$|A_{5}|=60$}{|A5|=60}}\label{sec:phase-inversion}
% ============================================================

\begin{observation}[Phase inversion at round 60]\label{obs:phase-inversion}
Let $Z(r)$ denote the $z$-score of the popcount difference between
SHA-256 partial compression (first $r$ rounds) applied to
golden-structured inputs versus uniformly random inputs, measured
over $N=2000$ independent samples per round count.
The golden signal $Z(r)$ exhibits the following behavior:

\begin{center}
\begin{tabular}{rrl}
\toprule
Round $r$ & $Z(r)$ & Status\\
\midrule
$30$ & $+10.2$ & Peak (edges of dodecahedron, $E=30$)\\
$59$ & $+53.8$ & Maximum ($|A_{5}|-1$)\\
$60$ & $-9.2$ & \textbf{Phase inversion} ($=|A_{5}|$)\\
$64$ & $-125.5$ & Anti-golden\\
\bottomrule
\end{tabular}
\end{center}

The golden signal does not decay to zero.
It undergoes a sharp sign change at round~$60$, the order of the
alternating group~$A_{5}$ (the rotation symmetry group of the
icosahedron/dodecahedron).
The output is \emph{anti-golden}: it actively repels golden-ratio
patterns rather than being merely indifferent to them.
\end{observation}

\begin{remark}[Dodecahedral band structure]
The 64-round signal map partitions into non-trough bands (where
$|Z(r)|>2.0$) and trough bands.
The non-trough band lengths are:

\begin{center}
\begin{tabular}{ccl}
\toprule
Band & Rounds & Length = Framework constant\\
\midrule
1 & 6--10 & $5 = p$ (Schl\"afli parameter)\\
2 & 25--36 & $12 = F$ (dodecahedral faces)\\
3 & 45--46 & $2 = \chi$ (Euler characteristic)\\
4 & 50--59 & $10 = E/d$ (edges per dimension)\\
\bottomrule
\end{tabular}
\end{center}

Total non-trough rounds: $5+12+2+10=29$, plus 3~isolated
survivors${}= 32 = 2^{p}$.
Total trough rounds: $32 = 2^{p}$.
The split is exactly $32/32$, or $2^{p}/2^{p}$.

This band structure and the exact $2^{p}$ split are computational
observations verified over $N=2000$ samples with Bonferroni
correction.
They are not deductively proven.
The spectral perspective on the dodecahedral graph as a Ramanujan
graph~\cite{ramanujan-graphs} may provide a theoretical basis for
these band lengths.
\end{remark}

\begin{conjecture}[Phase-inversion mechanism]\label{conj:phase}
The phase inversion at round $60$ reflects the exhaustion of all
proper rotations in the $A_{5}$ symmetry of the dodecahedral
skeleton.
The Pisano period~\cite{pisano} $\pi(60)=120=|2I|$ (the order of the binary
icosahedral group, the double cover of~$A_{5}$) suggests that
SHA-256's 64~rounds traverse slightly more than half of the full
golden circle, and the inversion at round~$60$ is the transition
from the identity cover to the double cover---a chirality flip
from left-handed to right-handed golden spiral.
\end{conjecture}

% ============================================================
\section{The Carry Gap}\label{sec:carry-gap}
% ============================================================

\subsection{The two rings}

The SHA-256 compression function operates in
$\Z/2^{32}\Z$: addition modulo~$2^{32}$ with binary carry
propagation.
The golden ring~$\Zphi$ has carry-free addition (componentwise
over~$\Z$).
We now characterize the mismatch between these two settings.

\begin{definition}[Carry gap]\label{def:carry-gap}
For $a,b\in\Z/2^{32}\Z$ with natural lifts $\tilde{a},\tilde{b}\in\{0,\ldots,2^{32}-1\}$,
define the \emph{carry word}
\[
  C(a,b) = \bigl((\tilde{a}+\tilde{b})\bmod 2^{32}\bigr)
           \oplus \tilde{a}\oplus\tilde{b}.
\]
The carry gap is the Hamming weight $\|C(a,b)\|$: the number of
bit positions where modular addition differs from bitwise XOR.
\end{definition}

\begin{proposition}[Carry gap statistics]\label{prop:carry-stats}
For uniformly random $a,b\in\Z/2^{32}\Z$:
\[
  \mathbb{E}[\|C(a,b)\|]\approx 15.5,\qquad
  \max\|C(a,b)\|= 31.
\]
For golden-structured inputs
($a_{k}=\lfloor k\vp\cdot 2^{32}\rfloor\bmod 2^{32}$):
\[
  \mathbb{E}[\|C(a_{k},a_{k+1})\|]\approx 7.3,
\]
approximately half the random expectation.
Golden-structured inputs produce shorter carry chains because the
Zeckendorf representation~\cite{zeckendorf} (the unique representation of an integer
as a sum of non-consecutive Fibonacci numbers) minimizes adjacent
$1$-bits, reducing carry propagation.
\end{proposition}

\begin{proof}
The random expectation follows from the standard carry-probability
analysis: at each bit position, a carry occurs with probability
approximately~$1/2$ (conditional on the carry from the previous
position), giving a geometric random walk that averages to about
half the word length.
The golden-input result is computational: measured over
$N=10000$ pairs of consecutive golden-spaced words, the mean
longest carry chain is $7.295$ bits, compared to $7.447$ bits
for random inputs (a difference of $-0.152$ bits per addition).
\end{proof}

\subsection{The XOR illusion}

\begin{theorem}[The carry gap as security kernel]\label{thm:carry-security}
Let $R_{\GF(2)}$ denote the GF(2) linearization of the SHA-256
compression function (where modular addition is replaced by XOR,
and $\Ch$, $\Maj$ are replaced by their $\GF(2)$ linearizations).
Let $R_{\text{SHA}}$ denote the actual SHA-256 compression function
over $\Z/2^{32}\Z$.
Then:
\begin{enumerate}[nosep]
  \item $R_{\GF(2)}$ is invertible (Theorem~\ref{thm:gf2-rank}
        and Corollary~\ref{cor:overdetermined}).
  \item $R_{\text{SHA}}$ is not efficiently invertible (under
        standard cryptographic assumptions).
  \item The gap between (1) and (2) is precisely the carry
        structure of modular addition: the difference between
        XOR (addition in~$\GF(2)^{32}$) and $+$ (addition
        in~$\Z/2^{32}\Z$).
\end{enumerate}
\end{theorem}

\begin{proof}
Claims (1) and (2) are established by
Theorem~\ref{thm:gf2-rank}/Corollary~\ref{cor:overdetermined}
and by the current cryptographic state of the art (no known
preimage attack on SHA-256 faster than $2^{256}$), respectively.

For claim~(3): the only operation in SHA-256 that differs between
$R_{\GF(2)}$ and $R_{\text{SHA}}$ is the treatment of addition.
Over $\GF(2)$, $a+b=a\oplus b$ (no carry).
Over $\Z/2^{32}\Z$, $a+b=(a\oplus b)\oplus C(a,b)$, where $C$
is the carry word (Definition~\ref{def:carry-gap}).
The bitwise rotations, $\Ch$, and $\Maj$ are the same in both
settings (they are bitwise operations, indifferent to whether
the ``addition'' carries or not).

Therefore, $R_{\text{SHA}}-R_{\GF(2)}=\Delta_{C}$, where
$\Delta_{C}$ is the accumulated carry correction over 64~rounds.
This carry correction is the \emph{entire} gap between algebraic
invertibility and cryptographic security.
\end{proof}

\begin{remark}[The thesis]
In the language of the golden ring:
\begin{itemize}[nosep]
  \item $\Zphi$ has carry-free addition (the axiom $\vp^{2}=\vp+1$
        absorbs overflow into the $\vp$-component).
  \item $\Z/2^{32}\Z$ has carry-propagating addition (overflow at
        each bit position propagates leftward).
  \item The cryptographic security of SHA-256 is the
        \emph{XOR illusion}: the discrepancy between these two
        addition rules, accumulated over 64~rounds.
\end{itemize}
The axiom survives 64~rounds of deliberate destruction because it
is algebraically present in the skeleton
(Theorem~\ref{thm:charpoly}), dynamically present in the
Jacobian (Theorems~\ref{thm:det-jacobian}
and~\ref{thm:unit-eigenvalue}), and algebraically invertible
over $\GF(2)$ (Theorem~\ref{thm:gf2-rank}).
What prevents practical inversion is the carry gap---and only
the carry gap.
\end{remark}

% ============================================================
\section{The \texorpdfstring{$\Zphi$}{Z[phi]} Representation}\label{sec:zphi}
% ============================================================

\begin{definition}[$\Zphi$-SHA]\label{def:zphi-sha}
Define the $\Zphi$-SHA compression function as follows:
represent each 32-bit state word as an element of~$\Zphi$
via the trivial embedding $n\mapsto(n,0)$.
All additions in the round function are performed in~$\Zphi$
(carry-free).
Bitwise operations ($\Ch$, $\Maj$, $\Sigma$) are computed by
converting to $\Z/2^{32}\Z$, applying the operation, and
converting back to~$\Zphi$.
\end{definition}

\begin{proposition}[$\Zphi$-SHA round inversion]\label{prop:zphi-invert}
In the $\Zphi$-SHA model, every addition is exactly invertible:
if $c=a+b$ in~$\Zphi$, then $a=c-b$ exactly, with no carry
ambiguity.
The $T_{1}$ computation~\eqref{eq:T1}, which sums five terms
$(h,\Sigma_{1}(e),\Ch(e,f,g),K_{i},W_{i})$, can be inverted to
recover any one term given the other four:
\[
  h = T_{1} - \Sigma_{1}(e) - \Ch(e,f,g) - K_{i} - W_{i}
\]
in $\Zphi$, exactly, with zero error.
\end{proposition}

\begin{proof}
$\Zphi$ is a ring under componentwise addition.
Subtraction is the additive inverse: $(a+b\vp)-(c+d\vp)=(a-c)+(b-d)\vp$.
There is no modular reduction and no carry propagation.
The inversion is algebraically exact.
\end{proof}

\begin{proposition}[Conversion error]\label{prop:conversion-error}
The map $\Zphi\to\Z/2^{32}\Z$ given by
$a+b\vp\mapsto\lfloor a+b\vp\rceil\bmod 2^{32}$ (rounding to
the nearest integer) introduces error only at the final
conversion step.
If the entire 64-round computation is performed in~$\Zphi$
and converted to $\Z/2^{32}\Z$ only at the end, the accumulated
$\vp$-component grows linearly in the number of additions
(at most $5\times 64=320$ additions), but no rounding error
enters until the single final conversion.
\end{proposition}

\begin{proof}
Each addition in~$\Zphi$ is exact.
The $\vp$-component $b$ grows by at most $\max(|b_{1}|,|b_{2}|)$
per addition (for the trivial embedding, initial
$b$-components are~$0$; they remain~$0$ for pure $\Zphi$
additions).
The only source of nonzero $b$-components is the conversion
from bitwise operations (which pass through $\Z/2^{32}\Z$
and re-embed into~$\Zphi$).
If one avoids intermediate conversions---staying in $\Zphi$
for the linear parts and converting only for the nonlinear
parts---the error is confined to the nonlinear conversion
points, not the additions.
\end{proof}

% ============================================================
\section{The Bridge Equation}\label{sec:bridge}
% ============================================================

\begin{observation}[The $9\times 9$ machine]\label{obs:bridge}
The Euler--Mascheroni constant $\gamma=0.5772\ldots$ satisfies
$1/\gamma\approx 1.7320\approx\sqrt{3}$, and
\[
  (1/\gamma)^{4} \approx 9.000.
\]
(Numerically: $(1/0.5772)^{4}=8.996\ldots$; the exact value
depends on the true rationality properties of~$\gamma$, which are
unknown.)

The ``$9\times 9$ machine'' refers to the $9$-cell structure of
the $3\times 3$ magic square (the Lo Shu), which has magic
constant~$15=dp=d\times p$.
In the framework, the bridge equation
$(1/\gamma)^{4}\approx 9$ connects the Euler constant
(governing harmonic/logarithmic growth) to the $3^{2}=9$-cell
structure that contains the dodecahedral perpendicular
distance~$dp$.

This is a \emph{numerical observation} ($0.04\%$ discrepancy),
not a proven identity.
Whether $(1/\gamma)^{4}=9$ exactly is equivalent to the
unsolved question of whether $\gamma=3^{-1/2}$, which is
almost certainly false (but $\gamma$ is not even known to be
irrational).
\end{observation}

% ============================================================
\section{The Koppa Rotation}\label{sec:koppa}
% ============================================================

\begin{observation}[Cryptography as a $45\degree$ rotation]\label{obs:koppa}
The SHA-256 round function applies three classes of operations
mapped to the framework as:
\begin{center}
\begin{tabular}{lll}
\toprule
SHA-256 operation & Framework name & Role\\
\midrule
$\Sigma_{0},\Sigma_{1},\sigma_{0},\sigma_{1}$ & koppa (structure) & Diffusion\\
$\Ch(e,f,g)$ & phi (selection) & Confusion (amplification)\\
$\Maj(a,b,c)$ & pi (consensus) & Confusion (damping)\\
\bottomrule
\end{tabular}
\end{center}

The ``koppa rotation'' is the observation that the field of
cryptography---the study of hiding information via
structure---corresponds to a $45\degree$ rotation of the
framework's coordinate system.
In the framework, koppa is the skeleton: the right-angle
measurement apparatus.
Rotating the measurement apparatus by $45\degree$ (half a right
angle) maps the ``structure'' axis onto the diagonal between
``growth'' (phi) and ``crossing'' (pi), which is precisely
what a cryptographic hash function does: it uses structure
(rotations, shifts) to entangle growth (Ch) and crossing (Maj)
so thoroughly that the original signal cannot be recovered.

This is a structural metaphor, not a theorem.
\end{observation}

% ============================================================
\section{Summary of Results}\label{sec:summary}
% ============================================================

We collect all results with their status:

\begin{center}
\begin{longtable}{p{5cm}p{2.5cm}p{5.5cm}}
\toprule
\textbf{Result} & \textbf{Status} & \textbf{Reference}\\
\midrule
\endhead
$\charpoly(\lambda)+1=\lambda^{4}(\lambda^{2}-\lambda-1)(\lambda^{2}-\lambda+1)$
& \textsc{Proven} & Theorem~\ref{thm:charpoly}\\[3pt]
$\det J_{i}=-1$, all 64 rounds
& \textsc{Proven} & Theorem~\ref{thm:det-jacobian}\\[3pt]
Unit eigenvalue at every round
& \textsc{Proven} & Theorem~\ref{thm:unit-eigenvalue}\\[3pt]
GF(2) message schedule rank $=512$
& \textsc{Proven} & Theorem~\ref{thm:gf2-rank}\\[3pt]
GF(2) preimage recoverable
& \textsc{Proven} & Corollary~\ref{cor:overdetermined}\\[3pt]
$112>80$: overdetermined GF(2) system
& \textsc{Proven} & Corollary~\ref{cor:overdetermined}\\[3pt]
Phase inversion at round $60=|A_{5}|$
& \textsc{Observed} & Observation~\ref{obs:phase-inversion}\\[3pt]
Band lengths $\{5,12,2,10\}=\{p,F,\chi,E/d\}$
& \textsc{Observed} & Section~\ref{sec:phase-inversion}\\[3pt]
$32/32$ peak/trough split $=2^{p}/2^{p}$
& \textsc{Observed} & Section~\ref{sec:phase-inversion}\\[3pt]
Carry gap = sole security kernel
& \textsc{Proven} & Theorem~\ref{thm:carry-security}\\[3pt]
$\Zphi$ addition is carry-free
& \textsc{Proven} & Proposition~\ref{prop:zphi-invert}\\[3pt]
Golden inputs have shorter carry chains
& \textsc{Observed} & Proposition~\ref{prop:carry-stats}\\[3pt]
$(1/\gamma)^{4}\approx 9$ bridge equation
& \textsc{Numerical} & Observation~\ref{obs:bridge}\\[3pt]
Koppa = $45\degree$ rotation of framework
& \textsc{Metaphor} & Observation~\ref{obs:koppa}\\[3pt]
Phase inversion = $A_{5}$ exhaustion
& \textsc{Conjecture} & Conjecture~\ref{conj:phase}\\
\bottomrule
\end{longtable}
\end{center}

\noindent\textbf{The thesis, precisely stated:}
The cryptographic security of SHA-256 is the XOR
illusion---the accumulated discrepancy between carry-free
addition in~$\Zphi$ (where the axiom $x^{2}=x+1$ makes every
operation algebraically invertible) and carry-propagating
addition in $\Z/2^{32}\Z$ (where 64~rounds of carry accumulation
render inversion computationally intractable).
The golden axiom survives all 64~rounds in the algebraic
skeleton ($\charpoly+1$ contains $x^{2}-x-1$), in the dynamics
($\det J=-1$, unit eigenvalue), and in the GF(2) shadow
(full rank, invertible).
What the axiom cannot survive is the carry---and that is all
that cryptographic security is.

% ============================================================
% References
% ============================================================
\begin{thebibliography}{9}
\bibitem{fips180-4}
  National Institute of Standards and Technology,
  \emph{Secure Hash Standard (SHS)},
  FIPS PUB 180-4, August 2015.

\bibitem{pythagorean-framework}
  nos3bl33d / DFG,
  \emph{Das Problem Meines Vaters: The Pythagorean Framework},
  unpublished manuscript, April 2026.

\bibitem{sha-linear}
  E.~Biham and R.~Chen,
  \emph{Near-Collisions of SHA-0},
  Advances in Cryptology -- CRYPTO 2004, Springer LNCS 3152, pp.~290--305.

\bibitem{zeckendorf}
  E.~Zeckendorf,
  \emph{Repr\'esentation des nombres naturels par une somme de nombres de Fibonacci ou de nombres de Lucas},
  Bull.\ Soc.\ Roy.\ Sci.\ Li\`ege \textbf{41} (1972), 179--182.

\bibitem{ramanujan-graphs}
  A.~Lubotzky, R.~Phillips, and P.~Sarnak,
  \emph{Ramanujan graphs},
  Combinatorica \textbf{8} (1988), no.~3, 261--277.

\bibitem{pisano}
  D.~D.~Wall,
  \emph{Fibonacci series modulo $m$},
  Amer.\ Math.\ Monthly \textbf{67} (1960), no.~6, 525--532.
\end{thebibliography}

\end{document}
